\documentclass[11pt]{article}

% packages
\usepackage{hyperref}
\usepackage{geometry}
\usepackage{enumitem}
\usepackage{booktabs}
\usepackage[backend=biber,style=numeric,sorting=none]{biblatex}
\addbibresource{references.bib}

% settings
\hypersetup{
    colorlinks=true,
    linkcolor=blue,
    urlcolor=blue,
}
\geometry{
    margin=1in,
}
\setlist[itemize]{leftmargin=1.3em}

\title{INCITE Report and Renewal Contribution}
\author{Qiu Group}
\date{}

\begin{document}

\maketitle
\tableofcontents

\section{Project Achievements}

\subsection{Significant Accomplishments to Date}

The INCITE allocation helped make advances in theoretical and computational
methods used to study self-trapped excitons (STEs) from first-principles
calculations. STEs have been shown to cause broadband white light emission in
certain double perovskites, showing coupling to phonons, Stokes shifts, and broad
emission peaks \cite{luo2018warmwhite,depaula2025time}. Being able to explain
this type of emission from first-principles calculations would help connect the
measured spectra to lattice distortions, phonon modes, and exciton-phonon
coupling parameters.

\subsubsection{Previous Methods and Challenges}

There are several ways to study STEs, but each method has
either computational challenges or methodological limitations.
The methods in the literature can be grouped
into three main approaches:
\begin{itemize}
\item Restricted Open-Shell Kohn-Sham method \cite{sahre2026roks}.
\item Excited-state forces \cite{ismailbeigi2003esf,delgrande2025esf}.
\item Minimization under the harmonic lattice approximation
    \cite{dai2024theory,dai2024firstprinciples}.
\end{itemize}

ROKS has been shown to describe STE states from density-functional theory (DFT)
\cite{sahre2026roks}. This can describe an excited-state structure by fixing
occupations, but it does not explain all effects that come from electron-hole
correlations in excitons.

Other methods account for electron-hole correlations by using excitons from the
Bethe-Salpeter Equation (BSE) under the GW approximation. Excited-state forces
(ESF) compute forces from these excitons and use them to relax the system into
the STE state \cite{ismailbeigi2003esf,delgrande2025esf}. ESF can capture the
large lattice distortion that forms an STE, including non-harmonic effects.
However, the calculation is expensive, especially in supercells, and it can
depend strongly on the initial lattice displacement used in the relaxation
process. A third approach is to minimize the exciton energy under a harmonic
lattice approximation \cite{dai2024theory,dai2024firstprinciples}. This is
computationally less expensive as it does not require a supercell, but it misses
non-harmonic lattice distortions.

One major limitation of all these methods is that they do not provide a way to
calculate the temperature dependence of STE properties, which has been observed
experimentally in \(\mathrm{Cs_2AgInCl_6}\) \cite{singh2024ultrafast}.

In this project, we made methodological developments and addressed
computational challenges based on these methods in the literature. The work
can be roughly grouped into two main points:
\begin{itemize}
\item Developing a many-body perturbation theory description of
STEs to explain temperature-dependent effects. This
includes approximating the energies using tadpole and Fan Migdal diagrams,
together with code development with partial self-consistency to obtain STE
states under the harmonic approximation.
\item Reducing the computational cost of ESF by using a better initial guess
from the self-energy method.
\end{itemize}

\subsubsection{Self-Energy Method Development}

We developed partially self-consistent self-energy formulations to describe
STE states, extending prior treatments of exciton-phonon coupling and
excitonic polarons \cite{antonius2022theory,dai2024theory,dai2024firstprinciples}.
The tadpole and Fan Migdal diagrams give a Dyson-like equation for the STE
states, which is then solved with partial self-consistency.

The tadpole diagram captures the static lattice distortion driven by the
exciton. The Fan Migdal term gives access to temperature-dependent effects that
minimization methods miss. Together, within the harmonic approximation, they
yield the phonon modes contributing to the STE emission peak and the
temperature dependence of the peak width.

Theory and code development consumed a large part of the allocation. Both
self-energy terms had to be implemented, debugged, and tested. Only the lowest
eigenvalues and eigenvectors of the self-energy matrix were computed, targeting
the low-energy STE states. Parallelization details are discussed in
Section~\ref{sec:performance}.

The method was first benchmarked on \(\mathrm{LiF}\) and \(\mathrm{SiO_2}\),
which have prior STE results in the literature and, for \(\mathrm{SiO_2}\),
experimental data. The quantities used for comparison included the STE energy,
the projection of the STE wavefunction onto exciton and phonon band states, and
the real-space localization of the wavefunction. These benchmarks identified the
phonon modes most responsible for STE formation and showed how the exciton
localizes in the STE state. The method was then applied to
\(\mathrm{Cs_2AgInCl_6}\). Compared to \(\mathrm{LiF}\),
\(\mathrm{Cs_2AgInCl_6}\) shows a stronger Fan Migdal contribution, driven by a
lower-energy optical mode that couples more strongly with the exciton: the
relevant phonon mode is around 15~meV in \(\mathrm{Cs_2AgInCl_6}\), compared to
around 60~meV in \(\mathrm{LiF}\).

Finite-momentum exciton calculations were carried out for all three materials, 
\(\mathrm{LiF}\), \(\mathrm{SiO_2}\), and \(\mathrm{Cs_2AgInCl_6}\), on \(2 \times 2 \times 2\), \(4 \times 4 \times 4\),
\(6 \times 6 \times 6\), and \(8 \times 8 \times 8\) grids, using Quantum
ESPRESSO for DFT, DFPT, and EPW calculations, and BerkeleyGW for GW and
GW-BSE calculations. These inputs were needed at each grid size
to extrapolate convergence toward an infinite grid limit. 

\subsubsection{Excited-State Forces and STE Relaxation}

The ESF method was introduced by Ismail-Beigi and Louie \cite{ismailbeigi2003esf}
and further developed by Del Grande and Strubbe \cite{delgrande2025esf}. ESF
calculations on \(\mathrm{LiF}\) and \(\mathrm{SiO_2}\) were used to validate
the workflow against prior literature results. For \(\mathrm{Cs_2AgInCl_6}\),
however, convergence difficulties
arose: some starting distortions led to slow convergence or oscillations. Since
each ESF step requires expensive electronic-structure work, repeatedly trying
different starting distortions is not practical. The self-energy calculation
gives a better initial guess, which reduced the number of ESF iterations needed
and made the calculation feasible on Frontier.

Each ESF iteration for \(\mathrm{Cs_2AgInCl_6}\) involved a full sequence of
DFT, DFPT, EPW, GW, and BSE calculations on a \(3\times3\times3\) supercell.
The initial lattice distortion for the first step was taken directly from the
self-energy calculation, ensuring a starting point tied to the STE rather than
an arbitrary displacement. At each subsequent step, the distortion is updated
from the forces computed in the previous iteration.

The quantities extracted from the ESF calculation are the total energy change
relative to the undistorted structure and the final electronic wavefunction,
which is compared against the exciton state to assess how much the STE state
resembles the original exciton and how it has changed due to the lattice
relaxation.

\subsubsection{Main Accomplishments and Impact}

The main work completed to date is:

\begin{itemize}
    \item A partially self-consistent Dyson-like equation with tadpole and Fan
        Migdal contributions was derived and implemented for predicting STE
        states under the harmonic approximation.
    \item The code was used to calculate STE states in \(\mathrm{LiF}\),
        \(\mathrm{SiO_2}\), and \(\mathrm{Cs_2AgInCl_6}\).
    \item The self-energy distortion was used as the initial guess for ESF
        calculations on \(\mathrm{Cs_2AgInCl_6}\), improving convergence
        and reducing computational cost.
    \item Properties including the Stokes shift, temperature dependence of the
        emission peak width, real-space STE wavefunction, and projection of the
        STE onto exciton and phonon components were computed to characterize
        STE behavior.
\end{itemize}

The methods developed here are general and can be applied to study STE states
in a wide range of materials, including two-dimensional systems and
Ruddlesden-Popper phases, where exciton-phonon coupling and self-trapping are
also of interest.

\subsection{Allocation Use}

\subsubsection{Usage to Date}

About 1,270 node hours have been used to date on Frontier by user
\texttt{krishnaa423}. Initially a lot of time was spent in
code development, debugging, and close analysis of intermediate results 
before running larger calculations. Table~\ref{tab:allocation}
shows the approximate breakdown by activity.

\begin{table}[h]
\centering
\small
\begin{tabular}{|l|l|l|}
\hline
\multicolumn{1}{|c|}{\textbf{Activity}} & \multicolumn{1}{c|}{\textbf{Node-Hour \%}} & \multicolumn{1}{c|}{\textbf{Systems}} \\
\hline
Self-energy code development, debugging, testing & $\sim$20\% &
    \(\mathrm{LiF}\) (primary test case) \\
\hline
Finite-momentum exciton input generation & $\sim$30\% &
    \(\mathrm{LiF}\), \(\mathrm{SiO_2}\), \(\mathrm{Cs_2AgInCl_6}\) \\
\hline
ESF iterations (DFT, DFPT, EPW, GW, BSE cycles) & $\sim$50\% &
    \(\mathrm{Cs_2AgInCl_6}\) \\
\hline
\end{tabular}
\caption{Approximate node-hour breakdown on Frontier.}
\label{tab:allocation}
\end{table}

\begin{itemize}
    \item \textbf{Self-energy code development, debugging, and testing
        ($\sim$20\%):} Initial implementation and validation of the self-energy
        code was done primarily on \(\mathrm{LiF}\), including debugging,
        benchmarking against literature results, and testing parallelization.
    \item \textbf{Finite-momentum exciton input generation ($\sim$30\%):} DFPT
        and EPW calculations via Quantum ESPRESSO, followed by GW-BSE
        calculations via BerkeleyGW, were run for
        \(\mathrm{LiF}\), \(\mathrm{SiO_2}\), and \(\mathrm{Cs_2AgInCl_6}\)
        on \(2\times2\times2\), \(4\times4\times4\), \(6\times6\times6\), and
        \(8\times8\times8\) grids. These inputs feed directly into the
        self-energy calculations.
    \item \textbf{ESF iterations ($\sim$50\%):} The dominant cost was the ESF
        calculation on \(\mathrm{Cs_2AgInCl_6}\), which required approximately
        six full DFT, DFPT, EPW, GW, and BSE cycles. Each iteration is
        expensive because it involves a complete electronic-structure and BSE
        calculation at each step.
\end{itemize}

\subsubsection{Projected Use}

For the remainder of the current allocation year, approximately 2,000 node
hours are anticipated for computing STE properties on the current systems, 
\(\mathrm{LiF}\), \(\mathrm{SiO_2}\), and \(\mathrm{Cs_2AgInCl_6}\), along
with associated code testing and debugging. This includes quantities such as
exciton projections, phonon contributions, real-space wavefunctions, and
temperature-dependent linewidths from both methods.
Following this, time will be spent on
manuscript preparation before pursuing calculations on new materials such as
two-dimensional systems or Ruddlesden-Popper phases. Usage will remain
intermittent rather than evenly spread.

\subsection{Application Parallel Performance}
\label{sec:performance}

The production work used established electronic-structure codes alongside the
in-house self-energy code. Quantum ESPRESSO and EPW handled DFT, DFPT, and
electron-phonon inputs. BerkeleyGW was used for GW and finite-momentum BSE
calculations. BerkeleyGW parallelizes well over Frontier: the finite-momentum
calculations over $\mathbf{Q}$ points were split into independent batches, and
the GW and BSE stages were parallelized over bands and $\mathbf{k}$ points.

The core computational step in the self-energy code is solving a Dyson-like
equation whose lowest-energy solutions give the STE eigenvalues and
eigenvectors. Rather than implement a custom parallel eigensolver, we used
\texttt{PETSc} and \texttt{SLEPc} through the Python interfaces
\texttt{petsc4py} and \texttt{slepc4py}, all compiled with \texttt{HIP} support
for Frontier's AMD GPUs. This provides distributed arrays and matrices across
GPUs and gives access to parallel eigensolvers that target only the lowest part
of the spectrum. This way the calculation is effectively parallelized over $Q$ points,
which made the \(4 \times 4 \times 4\), \(6 \times 6 \times 6\), and
\(8 \times 8 \times 8\) grid runs more practical. The resulting wavefunctions are
written to \texttt{.h5} files via \texttt{h5py} for parallel I/O.

The self-energy result also provided a better initial guess for the lattice
distortion used in the excited-state force calculations, reducing the number
of lattice displacement steps needed to reach the relaxed STE geometry.

The main remaining performance challenge is DFPT. A Quantum ESPRESSO build that
uses Frontier's AMD GPUs effectively for this part of the workflow has not been
identified, so input generation for the electron-phonon stages remains less
accelerated than the BerkeleyGW and self-energy stages.

\section{Project Plans for Next Year}

\subsection{Summary of Project Plans}

The next allocation period will focus on analyzing the completed calculations
and preparing manuscripts before starting any new calculations. The goal is to
extract STE observables from the data already in hand: the relaxed lattice
distortion, the Stokes shift, the emission linewidth, and the temperature
dependence of the photoluminescence spectrum.

The first priority is extracting and comparing properties from all completed
calculations across all systems and both methods. This means computing the STE
wavefunction, exciton projections onto band states, phonon-resolved
contributions, Stokes shifts, and temperature-dependent emission linewidths for
\(\mathrm{LiF}\), \(\mathrm{SiO_2}\), and \(\mathrm{Cs_2AgInCl_6}\), using
both the self-energy and ESF approaches wherever both are available. Comparing
the two methods across all three systems will show when the harmonic self-energy
picture is sufficient and when non-harmonic relaxation changes the result. The
goal at this stage is to have enough processed output to read and understand the
data, before any writing begins.

The second priority is manuscript preparation. Two papers are planned: one on
the self-energy framework for STEs and its temperature-dependent predictions,
and one on the \(\mathrm{Cs_2AgInCl_6}\) STE using the self-energy-guided ESF
result.

Extension to new materials such as two-dimensional systems or Ruddlesden-Popper
phases will be considered only after the current results are finalized and
the manuscripts are underway. Based on current usage, a new material requiring
finite-momentum exciton inputs and self-energy calculations is estimated at
roughly 1,000~node hours. Adding ESF iterations would bring the total to
roughly 2,000~node hours per material. These estimates could increase
significantly depending on the system size, since larger unit cells raise the
cost of the GW, BSE, and DFPT stages, and two-dimensional systems in particular
may require methodological development to parallelize well over Frontier.

\subsection{Development Work}
\label{sec:development}

Development during the current allocation period focused on the theory and
implementation of the STE self-energy method, including the self-consistent
tadpole term, the Fan Migdal term, finite-momentum exciton inputs, and
post-processing for STE observables. Development also included tests on
\(\mathrm{LiF}\), debugging, extension to denser grids, and the Frontier build
of \texttt{PETSc}, \texttt{SLEPc}, \texttt{petsc4py}, and
\texttt{slepc4py} with \texttt{HIP} support for AMD GPUs.

Future development work will focus on the post-processing code needed to
extract STE properties from completed calculations, including exciton
projections, phonon-resolved contributions, real-space wavefunctions, and
temperature-dependent linewidths. Some additional parallelization optimization 
may also be
needed depending on new systems and higher grid sizes. 

\section{Publications Resulting from INCITE Awards}

There are no submitted or published papers yet from this INCITE work that
include the required INCITE/OLCF acknowledgement.
Two manuscripts are in preparation from the calculations completed so far:

\begin{itemize}
    \item A paper on the first-principles exciton self-energy framework for
        STEs, including temperature-dependent effects and
        photoluminescence linewidths.
    \item A paper on \(\mathrm{Cs_2AgInCl_6}\), using self-energy-guided
        excited-state forces to obtain the STE distortion, Stokes shift, and
        emission properties of the double perovskite.
\end{itemize}

These papers will include the required INCITE/OLCF acknowledgements
when submitted.

\printbibliography

\end{document}
